%=========================================================================
\chapter{Introspection of WIR}
\label{chap:debug}
\renewcommand{\sectiontitle}{}
%=========================================================================

The classes, data structures and algorithms of the \gls{WIR} library are complex so that users might wish to get a deeper understanding of its inner workings without having to do a deep dive into its C++ code base and without having to rebuild \gls{WIR} completely.
Typical questions that come up when using \gls{WIR} are, e.g.: When are objects of a certain class created or destroyed? How does a complex analysis or optimization algorithm from \cref{chap:analysis,chap:optimization} behave in more detail? How does the control or data flow graph look like?

\gls{WIR} comes with a convenient mechanism allowing to do such kinds of introspection of the library's behavior without having to change the code base.
These mechanisms are called \emph{debug macros} in the following.
By using a simple configuration file, the desired output can selectively be turned on or off, giving the user full control over the amount of created introspection output.

The entire debug macro infrastructure gets activated only if a \gls{WIR} build is configured with option \texttt{-{}-enable-wir- debugmacros} (see \cref{tab:configure} on \cpageref{tab:configure}).
By default, the whole infrastructure is deactivated since it has a noticeable performance impact.
The rest of this chapter thus assumes a \gls{WIR} build configured with option \texttt{-{}-enable-wir-debugmacros}.
\begin{wirNote}
For release quality builds of the \gls{WIR} software package at the highest optimization level of the host compiler, users typically should not activate the debug macros.
\end{wirNote}

\cref{sec:debugmacros} first describes the general infrastructure provided by the debug macros.
After that, \cref{sec:wirmacros} presents how to use the debug macros integrated into the \gls{WIR} code for an in-depth introspection.


%-------------------------------------------------------------------------
\section{Debug Macro Framework}
\label{sec:debugmacros}
%-------------------------------------------------------------------------

The term debug macros originates from a couple of C++ preprocessor macros that are all declared in file \texttt{WIR/ LIBUSEFUL/libuseful/debugmacros.h}.
The entire framework consists of two parts:
\begin{itemize}
  \item the consequent and pervasive use of the debug macros in the C++ code of \gls{WIR} by the developers, and
  \item the dynamic and flexible configuration of actual introspection output to be produced for a \gls{WIR} user based on configuration file \texttt{wir.conf} that by default resides in directory \texttt{WIR/DEBUGMACRO\_CONF} (see configuration option \texttt{-{}-with-debugmacro-configuration} from \cref{tab:configure}).
\end{itemize}

The following discussion and the excerpt from \texttt{debugmacros.h} shown in \cref{lst:debugmacros} focuses on the most frequently used debug macros only which are heavily used by \gls{WIR}.
For a full description of all available debug macros, the interested reader is referred to header file \texttt{WIR/LIBUSEFUL/libuseful/debugmacros.h}.

\begin{lstlisting}[style=code,float=ht,caption={Most Relevant Debug Macros from \texttt{WIR/LIBUSEFUL/libuseful/debugmacros.h}},label=lst:debugmacros]
// Start debugging for the currently active C++ scope identified by the given activation name.
^#define DSTART^( ACTIVATIONNAME ) [...]Ü\label{lst:debugmacros1}Ü

// Produce debug output for the currently active C++ scope.
^#define DOUT^( ... ) [...]Ü\label{lst:debugmacros2}Ü

// Execute any arbitrary debugging code for the currently active C++ scope.
^#define DACTION^( ACTIONCODE ) [...]Ü\label{lst:debugmacros3}Ü
\end{lstlisting}

The most important debug macro is \texttt{DINIT} shown in \cref{lst:debugmacros1}.
It expects a string literal as argument which serves as so-called activation name.
The purpose of \texttt{DSTART} is to selectively activate the debugging of a complete C++ scope if the user has specified this in configuration file \texttt{wir.conf}.
Here, a C++ scope is any region in the C++ code base of \gls{WIR} between curly brackets \texttt{\{\ ...\ \}}.
For such a scope featuring \texttt{DSTART}, a message is printed out whenever this scope is entered during the execution of a user's \gls{WIR}-based application.
Likewise, an exit message is printed whenever such a scope is left, i.e., the user's application reaches the scope's closing \texttt{\}} bracket.

In order to give the user control for which scopes such enter and exit messages shall be printed, the activation name is used which serves as a unique identifier for debugging-relevant scopes.
The string literals used as activation names for \texttt{DSTART} all-over \gls{WIR} are hard-coded in the C++ code of \gls{WIR}.
By putting the activation name strings of the user's interest into the abovementioned configuration file \texttt{wir.conf}, debugging of the so-specified scope is activated and enter and exit messages are created.
If a scope's activation name is not present in the user's configuration file, no debug output will be produced for a scope.

Internally, the debug macros dynamically organize activation names in a stack.
Whenever a scope featuring \texttt{DSTART} is entered while executing a user's application, the scope's activation name is put on top of the current stack.
Similarly, the top of the stack is removed when leaving a \texttt{DSTART} scope.
This allows to properly keep track of the nesting of scopes while debugging \gls{WIR} applications that very frequently execute one scope from within another in a deeply nested fashion.

Besides such enter and exit messages, it regularly is desirable to print additional messages for debugging purposes.
This can include, e.g., simple messages (\texttt{"Hello world."}) indicating that a certain position in the code has been reached, or messages printing the current contents of important variables etc.
For this purpose, the \texttt{DOUT} macro (\cref{lst:debugmacros2}) is used within \gls{WIR}.
It accepts a variadic number of arguments that are all treated as regular C++ code serving to produce output to a \texttt{std::ostream}.
This way, complex debug messages can be created as follows:
\begin{lstlisting}[style=code]
DOUT( "Current address: 0x" æ<<æ ästd::hexä æ<<æ addr æ<<æ "." << ästd::decä æ<<æ ästd::endlä );
\end{lstlisting}
\texttt{DOUT} only generates actual output if it is part of a scope featuring \texttt{DSTART} and if the activation name provided to \texttt{DSTART} shows up in the user's configuration file.

Finally, it might be necessary to execute complex code for debugging purposes instead of printing simple messages.
This usually is the case when data structures shall be inspected and checked for sanity.
This kind of debugging is supported by macro \texttt{DACTION} from \cref{lst:debugmacros3}.
It expects any kind of C++ code to be executed only if, once again, the activation name of the scope is configured by the user.
The following example displays the names of all childs of some register \texttt{reg} only if debugging of the scope containing this \texttt{DACTION} is desired by the user:
\begin{lstlisting}[style=code]
DACTION(
  for ( WIR_VirtualRegister æ&æchild : reg.getChilds() )
    DOUT( "Current child register: " æ<<æ child.getName() æ<<æ ästd::endlä ); );
\end{lstlisting}

For \gls{WIR} builds configured with disabled debug macros, all preprocessor definitions from file \texttt{debugmacros.h} get replaced by empty dummies so that the manifold occurrences of \texttt{DSTART}, \texttt{DOUT} and \texttt{DACTION} within the \gls{WIR} library leave absolutely no trace in this case.

Debug macro configuration files are simple plain-text files that can be edited by the user in any text editor.
They are interpreted line by line, empty lines and comment lines starting with "\texttt{\#}" or "\texttt{//}" are ignored.
One line in a debug macro configuration file now simply contains exactly one activation name of some scope to be debugged.
The strings used in the context of \texttt{DSTART} and within a user's configuration file are always case-sensitive, and they are interpreted literally, i.e., spacing or punctuation also matters.
Thus, the configuration file needs to list the exact and identical strings as they show up in the \texttt{DSTART} directives to be debugged.


%-------------------------------------------------------------------------
\section{WIR Debug Macros}
\label{sec:wirmacros}
%-------------------------------------------------------------------------

As mentioned earlier, \gls{WIR} by default considers the debug macro configuration file \texttt{WIR/DEBUGMACRO\_CONF/wir. conf}.
For the user's convenience, a configuration template listing all currently existing and supported activation names is provided in file \texttt{WIR/DEBUGMACRO\_CONF/wir.conf.bak}.
Activation names are all grouped by their associated \gls{WIR} classes.
This template gives an overview for which parts of \gls{WIR} debugging can be activated.
By copying and pasting strings from the template to \texttt{wir.conf}, an appropriate configuration can be crated easily.

As can be seen, these activation names regularly refer to the various methods of \gls{WIR} classes, including their argument and return types.
This results from the fact that within the entire C++ code base of the \gls{WIR} library, all methods always consequently start with a dedicated \texttt{DSTART} directive.
Like this, the begin and end of execution of individual \gls{WIR} methods can be traced by the user.

\cref{lst:variadicconstruction} on \cpageref{lst:variadicconstruction} showed an example for the variadic construction of object hierarchies within \gls{WIR}.
When using the following debug macro configuration:
\begin{lstlisting}[style=code]
WIR_CompilationUnit::WIR_CompilationUnit(WIR_Function&&, Args&& ...)
WIR_Function::WIR_Function(string&&)
WIR_BasicBlock& WIR_Function::pushBackBasicBlock(WIR_BasicBlock&&)
WIR_BasicBlock::WIR_BasicBlock(WIR_Instruction&&, Args&& ...)
WIR_Instruction::WIR_Instruction(WIR_Operation&&, Args&& ...)
WIR_Operation::WIR_Operation(ÜconstÜ WIR_BaseProcessor::OpCode&, ÜconstÜ WIR_BaseProcessor::OperationFormat&, Args&& ...)
\end{lstlisting}
the creation of the various objects by the example program can be observed.
Compiling and executing the example from \cref{lst:variadicconstruction} with the above configuration file in place produces the following output:
\begin{lstlisting}[style=code]
^--> enter^ WIR_Function::WIR_Function(string&&) (0)
Ü\color{red}<-{}-  exitÜ WIR_Function::WIR_Function(string&&) (0)
^--> enter^ WIR_CompilationUnit::WIR_CompilationUnit(WIR_Function&&, Args&& ...) (0)
Ü\color{red}<-{}-  exitÜ WIR_CompilationUnit::WIR_CompilationUnit(WIR_Function&&, Args&& ...) (0)
^--> enter^ WIR_Operation::WIR_Operation(ÜconstÜ WIR_BaseProcessor::OpCode&, ÜconstÜ WIR_BaseProcessor::OperationFormat&, Args&& ...) (0)
Ü\color{red}<-{}-  exitÜ WIR_Operation::WIR_Operation(ÜconstÜ WIR_BaseProcessor::OpCode&, ÜconstÜ WIR_BaseProcessor::OperationFormat&, Args&& ...) (0)
^--> enter^ WIR_Instruction::WIR_Instruction(WIR_Operation&&, Args&& ...) (0)
Ü\color{red}<-{}-  exitÜ WIR_Instruction::WIR_Instruction(WIR_Operation&&, Args&& ...) (0)
^--> enter^ WIR_Operation::WIR_Operation(ÜconstÜ WIR_BaseProcessor::OpCode&, ÜconstÜ WIR_BaseProcessor::OperationFormat&, Args&& ...) (0)
Ü\color{red}<-{}-  exitÜ WIR_Operation::WIR_Operation(ÜconstÜ WIR_BaseProcessor::OpCode&, ÜconstÜ WIR_BaseProcessor::OperationFormat&, Args&& ...) (0)
^--> enter^ WIR_Instruction::WIR_Instruction(WIR_Operation&&, Args&& ...) (0)
Ü\color{red}<-{}-  exitÜ WIR_Instruction::WIR_Instruction(WIR_Operation&&, Args&& ...) (0)
^--> enter^ WIR_BasicBlock::WIR_BasicBlock(WIR_Instruction&&, Args&& ...) (0)
Ü\color{red}<-{}-  exitÜ WIR_BasicBlock::WIR_BasicBlock(WIR_Instruction&&, Args&& ...) (0)
^--> enter^ WIR_BasicBlock& WIR_Function::pushBackBasicBlock(WIR_BasicBlock&&) (0)
Ü\color{red}<-{}-  exitÜ WIR_BasicBlock& WIR_Function::pushBackBasicBlock(WIR_BasicBlock&&) (0)
\end{lstlisting}

Changing the previous debug macro configuration to:
\begin{lstlisting}[style=code]
ÜvoidÜ WIR_System::createComponents(ÜconstÜ string&)
\end{lstlisting}
and re-executing the very same test program now yields the following output:
\begin{lstlisting}[style=code]
^--> enter^ ÜvoidÜ WIR_System::createComponents(ÜconstÜ string&) (1)
  Reading system configuration file '/tmp/WIR/arch/generic/genericmips.sys'.
  Created ÜnewÜ processor core 'CORE0'
    Instruction set architecture: 'WIR generic'
    Clock frequency: 150000000 Hz
    Voltage: 3.3 V
  Created ÜnewÜ memory region 'RAM'
    Base address: 0x0
    Length: 0x100000
    Attributes: 0xf
    Delay: 1
    Memory hierarchy: { 'CORE0' }
Ü\color{red}<-{}-  exitÜ ÜvoidÜ WIR_System::createComponents(ÜconstÜ string&) (1)
\end{lstlisting}
As can be seen from this example, just editing the plain-text configuration file helps to display completely different output, now showing the details how a \gls{WIR} system is set up for the generic MIPS architecture from \cref{ssec:mipsmodel}.

Besides activation names referring to complete methods of the \gls{WIR} code base, a few particular ones are provided for the introspection of complex analyses and optimizations.
They are summarized in \cref{lst:advancedmacros}.
The control tree iteratively constructed during structural control flow analysis (see \cref{sec:structcflow} at \cpageref{sec:structcflow}) can be visualized using the activation names from \cref{lst:advancedmacros1,lst:advancedmacros2}.
The former shows the current control tree after each individual iteration of the analysis, while the latter shows only the finally created result.

\cref{lst:advancedmacros3,lst:advancedmacros4} show the activation names relating to the visualization of the bit-true \gls{DFG} from \cref{def:bitdfg} on \cpageref{def:bitdfg}.
The first one enables the graph visualization right after the top-down analysis, i.e., the forward data flow analysis.
The final bit-true \gls{DFG} resulting from both top-down and bottom-up backward analysis can be shown using the activation name from \cref{lst:advancedmacros4}.

Due to its inherent complexity, various debugging facilities are provided for the graph coloring-based register allocation described in \cref{sec:optregalloc} on \cpageref{sec:optregalloc}.
Most of them serve to visualize the interference graph during the various steps of the graph coloring algorithm: immediately after having built the initial, uncolored interference graph (\cref{lst:advancedmacros5}); after graph simplification, once showing only the remaining non-simple nodes (\cref{lst:advancedmacros6}), or once showing both simple and non-simple nodes (\cref{lst:advancedmacros7}); after coalescing of graph nodes (\cref{lst:advancedmacros8}); after coloring the graph (\cref{lst:advancedmacros9}); and finally when applying stack coalescing (\cref{lst:advancedmacros10}).

\begin{lstlisting}[style=code,float=ht,caption={Activation Names for the In-Depth Introspection of WIR Analyses and Optimizations},label=lst:advancedmacros]
Üvirtual voidÜ WIR_StructuralAnalysis::runAnalysis(WIR_Function&).visualizeÜ\label{lst:advancedmacros1}Ü
Üvirtual voidÜ WIR_StructuralAnalysis::runAnalysis(WIR_Function&).ÜfinalÜ.visualizeÜ\label{lst:advancedmacros2}Ü

Üvirtual voidÜ WIR_BitDFA::runAnalysis(WIR_Function&).topdown.visualizeÜ\label{lst:advancedmacros3}Ü
Üvirtual voidÜ WIR_BitDFA::runAnalysis(WIR_Function&).bottomup.visualizeÜ\label{lst:advancedmacros4}Ü

ÜvoidÜ WIR_GraphColoring::build(WIR_Function&, WIR_InterferenceGraph&).visualizeÜ\label{lst:advancedmacros5}Ü
Üvirtual voidÜ WIR_GraphColoring::runOptimization(WIR_Function&).simplified.visualizeÜ\label{lst:advancedmacros6}Ü
Üvirtual voidÜ WIR_GraphColoring::runOptimization(WIR_Function&).simplified.visualizeFullÜ\label{lst:advancedmacros7}Ü
WIR_BaseRegister& WIR_GraphColoring::combine(ÜconstÜ WIR_BaseRegister&, ÜconstÜ WIR_BaseRegister&, WIR_InterferenceGraph&).visualizeÜ\label{lst:advancedmacros8}Ü
Üvirtual voidÜ WIR_GraphColoring::runOptimization(WIR_Function&).colored.visualizeFullÜ\label{lst:advancedmacros9}Ü
ÜvoidÜ WIR_GraphColoring::computeStackLocations(ÜconstÜ WIR_Function&).visualizeÜ\label{lst:advancedmacros10}Ü
WIR_GraphColoring.worklistsÜ\label{lst:advancedmacros11}Ü
WIR_GraphColoring.invariantsÜ\label{lst:advancedmacros12}Ü

WIR_JumpCorrection.invariantsÜ\label{lst:advancedmacros13}Ü
TC_JumpCorrection.invariantsÜ\label{lst:advancedmacros14}Ü

ÜvoidÜ WIR_SchedulingRegion::buildDG().visualizeÜ\label{lst:advancedmacros15}Ü
\end{lstlisting}

According to~\cite[chapter~16.3]{Much97}, the \gls{WIR} graph coloring is implemented as a typical worklist algorithm.
The processing of the various worklists, i.e., when operations or registers are added to or removed from certain worklists can be monitored by configuring the activation name from \cref{lst:advancedmacros11}.
Furthermore, \cite[chapter~16.3]{Much97} and~\cite[chapter~11]{Appel04} formulate a couple of invariants, i.e., structural properties of the interference graph and/or the worklists that must always hold.
The explicit verification of all these invariants (and many others, going beyond the textbooks cited above) and thus in-depth sanity checks of all algorithms can be turned on using the activation name from \cref{lst:advancedmacros12}.
However, the computational overhead of this introspection of worklists and invariants is extremely high.
Even if \cref{lst:advancedmacros11,lst:advancedmacros12} are not active, the very frequent executions of \texttt{DSTART} and the corresponding comparisons of activation names with configuration file contents still lead to a noticeable performance degradation.
For this reason and in order to obtain a reasonable performance of register allocation even for \gls{WIR} builds configured with \texttt{-{}-enable-wir-debugmacros}, these two activation names are ``mechanically'' turned off completely by default.
Should a \gls{WIR} user wish to have a closer look at worklists or invariants, it thus is not sufficient to just copy \cref{lst:advancedmacros11,lst:advancedmacros12} into configuration file \texttt{wir.conf}.
Instead, the definition of macro \texttt{ENABLE\_INVARIANTS} at the top of file \texttt{WIR/optimizations/gcregallocation/wirgraphcoloring.cc} needs to be activated and the \gls{WIR} library rebuilt once afterwards.

The very same applies to invariants that check the memory layout and start addresses of basic blocks during jump correction (see \cref{sec:optjumps} on \cpageref{sec:optjumps}).
This checking of invariants is in general supported by the activation names from \cref{lst:advancedmacros13,lst:advancedmacros14}.
But again and for the same reasons as discussed in the previous paragraph, the macro \texttt{ENABLE\_INVARIANTS} must be defined in files \texttt{WIR/optimizations/jumpcorrection/wirjumpcor\-rec\-tion.cc} and \texttt{WIR/arch/tricore/optimizations/jumpcorrection/tcjumpcorrection.cc}, resp.

Last but not least, the activation name from \cref{lst:advancedmacros15} serves to visualize the dependence graph showing \gls{RAW}, \gls{WAR} and many other dependencies between \gls{WIR} operations during instruction scheduling, cf. \cref{sec:optsched} on \cpageref{sec:optsched}.


\cleardoubleoddpage
